\documentclass[10pt]{article}
\usepackage[margin=0.72in]{geometry}
\usepackage{amsmath,amssymb,amsthm,mathtools}
\usepackage{graphicx}
\usepackage{microtype}
\usepackage{xcolor}
\usepackage[colorlinks=true,linkcolor=blue!50!black,citecolor=blue!50!black,urlcolor=blue!50!black]{hyperref}

\newtheorem{theorem}{Theorem}
\newtheorem{lemma}{Lemma}
\newcommand{\A}{\mathcal A}
\newcommand{\LUC}{\operatorname{LUC}}
\newcommand{\ip}[2]{\left\langle #1,#2\right\rangle}
\setlength{\parindent}{0pt}
\setlength{\parskip}{4pt}

\title{An $F\ell^1$-basis need not give $\eta$-factorization in $\LUC(\A)$}
\author{Candidate counterexample for arXiv:2412.15029}
\date{27 August 2026}

\begin{document}
\maketitle

\vspace{-1.2em}
\begin{center}
\fbox{\parbox{0.94\textwidth}{\textbf{Status.}
Candidate full counterexample, likely valid, pending expert review.}}
\end{center}

\section*{Source question}

Filali and Galindo prove that an $F\ell^1(\eta)$-basis gives
$\eta$-factorization in $\A^*$ and one-factorization in $\LUC(\A)$.  They ask
in Remark~4.15(iii) whether it also gives $\eta$-factorization in
$\LUC(\A)$ \cite[p.~20]{FilaliGalindo}.

\begin{center}
\includegraphics[width=0.84\textwidth]{figures/open_problem_crop.png}
\end{center}
\vspace{-0.5em}
\textbf{Transcription.}  In the case $m=1$, does every bounded family
$\{S_\lambda:\lambda<\eta\}\subset\LUC(\A)$ admit
\[
 S_\lambda=\Psi_\lambda\cdot T\qquad(\lambda<\eta),
\]
where $(\Psi_\lambda)$ is bounded in $\LUC(\A)^*$ and one common
$T$ belongs to $\LUC(\A)$?

\section*{Counterexample}

Take the convolution algebra $\A=L^1(\mathbb R)$.  Its compact covering is
$\aleph_0$, and source Theorem~6.1 states that $L^1(\mathbb R)$ has an
$F\ell^1(\aleph_0)$-basis of type~2 \cite[p.~28]{FilaliGalindo}.  Under the
standard identification $\LUC(\A)=\LUC(\mathbb R)$, put
\[
 S_n(x)=e^{inx},\qquad n=1,2,\ldots.
\]
This is a bounded family in $\LUC(\A)$, with $\|S_n\|_\infty=1$.

\section*{Proof Intuition}

A common factor $T\in\LUC(\A)$ imposes a uniform continuity condition on all
its bounded left factors: a bounded approximate identity must converge on the
whole family uniformly.  High-frequency characters violate precisely this
condition.  Each approximate-identity element has a Fourier transform tending
to zero at high frequency, so it moves some character by almost its full norm.

\begin{lemma}[Uniform-essentiality obstruction]\label{lem:uniform}
Let $\A$ have a bounded right approximate identity $(u_j)$.  If a bounded
family in $\LUC(\A)$ has a factorization
\[
 S_\lambda=\Psi_\lambda\cdot T,
 \qquad \sup_\lambda\|\Psi_\lambda\|<\infty,
 \qquad T\in\LUC(\A),
\]
then
\[
 \sup_\lambda\|S_\lambda\cdot u_j-S_\lambda\|\longrightarrow0.
\]
\end{lemma}

\begin{proof}
First, $T\cdot u_j\to T$ in norm.  Indeed, by definition
$\LUC(\A)=\overline{\operatorname{span}}(\A^*\cdot\A)$.  For a generator
$R\cdot a$,
\[
 (R\cdot a)\cdot u_j=R\cdot(au_j)\longrightarrow R\cdot a,
\]
and boundedness of $(u_j)$ extends this convergence to the closed span.
Using the module associativity recorded in the source,
\[
 (\Psi_\lambda\cdot T)\cdot u_j
 =\Psi_\lambda\cdot(T\cdot u_j).
\]
Consequently, if $C=\sup_\lambda\|\Psi_\lambda\|$, then
\[
 \sup_\lambda\|S_\lambda\cdot u_j-S_\lambda\|
 \le C\|T\cdot u_j-T\|\longrightarrow0. \qedhere
\]
\end{proof}

\begin{theorem}\label{thm:counterexample}
$L^1(\mathbb R)$ has an $F\ell^1(\aleph_0)$-basis, but
$\LUC(L^1(\mathbb R))$ does not have the $\aleph_0$-factorization property.
Thus the implication asked for in Remark~4.15(iii) is false.
\end{theorem}

\begin{proof}
Let $(u_j)$ be any bounded approximate identity in $L^1(\mathbb R)$.  With
the usual dual pairing and convolution product, direct calculation gives
\begin{align*}
 \ip{S_n\cdot u_j}{b}
 &=\ip{S_n}{u_j*b}\\
 &=\left(\int_{\mathbb R}u_j(t)e^{int}\,dt\right)
   \ip{S_n}{b}
 \qquad(b\in L^1(\mathbb R)).
\end{align*}
Hence
\[
 S_n\cdot u_j=m_j(n)S_n,
 \qquad m_j(n)=\int_{\mathbb R}u_j(t)e^{int}\,dt.
\]
(The opposite pairing convention merely replaces $n$ by $-n$.)  For every
fixed $j$, the Riemann--Lebesgue lemma \cite[Ch.~1]{Rudin} yields
$m_j(n)\to0$ as $n\to\infty$.  It follows that
\[
 \sup_{n\ge1}\|S_n\cdot u_j-S_n\|_\infty
 =\sup_{n\ge1}|m_j(n)-1|\ge1
 \qquad\text{for every }j.
\]

If the requested factorization existed, there would be a bounded sequence
$(\Psi_n)\subset\LUC(\A)^*$ and a common $T\in\LUC(\A)$ with
$S_n=\Psi_n\cdot T$.  Lemma~\ref{lem:uniform} would make the preceding
supremum tend to zero, a contradiction.  The source hypothesis was verified
above from Theorem~6.1, so this is the desired counterexample.
\end{proof}

\section*{Concrete Gaussian check}

For the normalized Gaussian approximate identity
\[
 u_j(t)=\frac{j}{\sqrt\pi}e^{-j^2t^2},
\]
one has $m_j(n)=e^{-n^2/(4j^2)}$.  Thus already at $n=20j$ the defect is
$1-e^{-100}$, independently of the scale $j$.  The supplied script checks
this numerical illustration; the proof above applies to every bounded
approximate identity.

\section*{Verification, scope, and novelty}

The source hypothesis, target-space identification, module associativity,
approximate-identity convergence, Fourier multiplier, and uniform contradiction
were audited separately.  The result does not affect source Theorem~4.16:
one-factorization can hold even though a bounded countable family has no common
factor.

A bounded search on 27 August 2026 checked the run indexes, the current official
arXiv record, exact phrases from Remark~4.15(iii), the paper title and authors,
and combinations of $F\ell^1$, $\LUC(\A)$, factorization, characters, and
approximate identities.  Only the source question was found, and the arXiv
record still lists v1.  Novelty is plausible but provisional.

\begin{thebibliography}{9}
\bibitem{FilaliGalindo}
M. Filali and J. Galindo,
\emph{$\ell^1$-bases, algebraic structure and strong Arens irregularity of
Banach algebras in harmonic analysis}, arXiv:2412.15029v1, 2024.

\bibitem{Rudin}
W. Rudin,
\emph{Fourier Analysis on Groups}, Interscience, New York, 1962.
\end{thebibliography}

\end{document}
